\mainchapterstyle
\chapter{روش تحقیق}
\pagestyle{mainheader} 


\section{مقدمه}
شرح کامل روش تحقیق است. این فصل بسته به نوع روش تحقیق و با نظر استاد راهنما می‌تواند مواد و روش‌ها\LTRfootnote{Materials and Methods}
نام گیرد. این فصل حدود 15 صفحه است.


\section{محتوا (نام گذاری بر اساس روش تحقیق و مسئله مورد مطالعه)}

\subsection{علت انتخاب روش}
دلیل یا دلایل انتخاب روش تحقیق را تشریح می‌کند.

\subsection{تشريح كامل روش تحقیق}
برای اینکه پایان‌نامه دارای ارزش علمی باشد، باید قابل تکرار باشد و داوران و خوانندگان از امکان
تکرارپذیر بودن کار شما مطمئن شوند. شما باید اساس تکرار آزمایش را به‌وسیله دیگران در این قسمت فراهم
کنید. تکرارپذیری آزمایشات و روش شما با میزان پتانسیل تکرار نتایج برابر یا نزدیک به آن است، در زیر به
تعدادی از روش‌های تحقیق اشاره شده است:

\begin{itemize}[leftmargin=*] % example of titled bullets

    \item \textbf{روش تحقیق مطالعه‌ی موردی} \\
    توصیف کامل برنامه‌ی آزمایشگاهی شامل مواد مصرفی و نحوه‌ی ساخت نمونه‌ها، شرح آزمایش‌ها شامل نحوه تنظیم و آماده‌سازی آزمایش‌ها و دستگاه‌های مورد استفاده و دقت و نحوه‌ی کالیبره کردن، شرح دستگاه ساخته شده (در صورت ساخت) و ارائه‌ی روش اعتبارسنجی.
    \item \textbf{روش تحقیق آماری} \\
    توصیف ابزارهای گردآوری اطلاعات کمی و کیفی، اندازه‌ی نمونه‌ها، روش نمونه‌برداری، تشریح مبانی روش آماری و ارائه‌ی روش اعتبارسنجی.
    \item \textbf{روش تحقیق نرم‌افزارنويسی} \\
    توصیف کامل برنامه‌نویسی، مبانی برنامه و ارائه‌ی روش اعتبارسنجی.

    \item \textbf{روش تحقیق مطالعه‌ی موردی} \\
    توصیف کامل محل و موضوع مطالعه، علت انتخاب مورد و پارامترهایی که تحت ارزیابی قرار داده می‌شوند، و ارائه‌ی روش اعتبارسنجی.
    \item \textbf{روش تحقیق تحلیلی يا مدل‌سازی} \\
    توصیف کامل مبانی یا اصول تحلیل یا مدل و ارائه‌ی روش اعتبارسنجی. در ارائه مدل ریاضی معمولا نیاز است اندیس ها، پارامترها، متغییرهای تصمیم و فرمول‌های مدل به صورت سیستماتیک ارائه شود. پیشنهاد می‌گردد برای نمایش اندیس ها، پارامتر‌ها و متغییرهای تصمیم از سه جدول به صورت زیر استفاده گردد.
    
    \begin{table}[h] % example of using tables
    \centering
    \caption{نمونه‌ی استفاده از جدول}
    \label{table:example}
    \begin{tabular}{cc}
    \hline
    ستون اول & ستون دوم \\
    \hline
    ۱ & ۲ \\
    \hline
    \end{tabular}
    \end{table}

    جدول~\ref{table:example} نمونه‌ای از استفاده از جدول را نشان می‌دهد در جدول \ref{longtable:example} هم میتوانید یک نمونه از جدول طولانی را ببینید


    \begin{longtable}{|c|c|c|} % example of using long table
    \caption{نمونه‌ی استفاده از جدول طولانی} \\
    \label{longtable:example}
    \hline
    ستون ۱ & ستون ۲ & ستون ۳ \\
    \hline
    \endfirsthead

    % ---------- continued title ----------
    \caption*{ادامه جدول \thetable} \\
    \hline
    ستون ۱ & ستون ۲ & ستون ۳ \\
    \hline
    \endhead

    % ---------- table body ----------
    داده & داده & داده \\
    \hline
    داده & داده & داده \\
    \hline
    داده & داده & داده \\
    \hline
    داده & داده & داده \\
    \hline
    داده & داده & داده \\
    \hline
    داده & داده & داده \\
    \hline
    داده & داده & داده \\
    \hline
    داده & داده & داده \\
    \hline
    داده & داده & داده \\
    \hline
    داده & داده & داده \\
    \hline
    داده & داده & داده \\
    \hline
    داده & داده & داده \\
    \hline
    داده & داده & داده \\
    \hline
    داده & داده & داده \\
    \hline
    داده & داده & داده \\
    \hline
    داده & داده & داده \\
    \hline
    داده & داده & داده \\
    \hline
    داده & داده & داده \\
    \hline
    داده & داده & داده \\
    \hline
    داده & داده & داده \\
    \hline
    داده & داده & داده \\
    \hline
    داده & داده & داده \\
    \hline
    داده & داده & داده \\
    \hline
    داده & داده & داده \\
    \hline
    داده & داده & داده \\
    \hline
    داده & داده & داده \\
    \hline
    داده & داده & داده \\
    \hline
    داده & داده & داده \\
    \hline
    داده & داده & داده \\
    \hline
    داده & داده & داده \\
    \hline
    داده & داده & داده \\
    \hline
    داده & داده & داده \\
    \hline
    داده & داده & داده \\
    \hline
    داده & داده & داده \\
    \hline
    داده & داده & داده \\
    \hline
    داده & داده & داده \\
    \hline
    داده & داده & داده \\
    \hline
    داده & داده & داده \\
    \hline
    داده & داده & داده \\
    \hline

    % ... many rows ...
    \end{longtable}

    \(\min(\delta_1, \delta)\)  نمونه‌ای از استفاده \lr{inline math} را نشان می‌دهد در حالی که در \ref{eq:sample} میتوانید نمونه‌ای از عبارات ریاضی را ببینید

    \begin{equation}
        \label{eq:sample}
        \bar{x} = \frac{\sum_{i=1}^{n} x_i}{n}
    \end{equation}


    همچنین در \ref{align:sample} خواهیم داشت:
    \begin{align}
    \label{align:sample}
    \mathbb{E}[Z] 
    &= \mathbb{E}[Z_{>\delta}] + \mathbb{E}[Z_{<\delta}] \\
    &\geq |I^{>\delta}| \cdot E_1 + |I^{<\delta}| \cdot \left(\frac{1}{2} - \frac{\mu\beta}{2}\right) \\
    &\geq p|y| \left[\frac{1}{2} + \left(\delta - \frac{1}{2}\right)\mu\beta\right] 
    + (1-p)|y| \left(\frac{1}{2} - \frac{\mu\beta}{2}\right) \\
    &= |y| \left[ \frac{1}{2} + p\delta\mu\beta - \frac{\mu\beta}{2} \right] \\
    &= \frac{|y|}{2} + \left(p\delta - \frac{1}{2}\right) \mu\beta |y|.
    \end{align}

    \item \textbf{روش تحقیق میدانی}\\
    چگونگی دستیابی به داده‌ها در میدان عمل و نحوه برداشت از پاسخ‌های دریافتی.

\end{itemize}



\begin{algorithm}[H] % example of using an algorithm
    \caption{\(\mathsf{Insertion Sort}(A)\)}
    \label{alg:inssort}
    \begin{algorithmic}[1]
        \REQUIRE آرایه \(A\)
        \ENSURE آرایه \(A\) در حالت مرتب شده
        \FOR{\(j = 2\) \TO \(A.length\)}
            \STATE \(key \gets A[j]\)
            \STATE \(i \gets j - 1\) \hfill // عنصر \(A[j]\) را در دنباله مرتب شده \(A[1,\ldots,j-1]\) درج کنید
            \WHILE{\(i > 0 \) \AND \(A[i] > key \)}
                \STATE \(A[i + 1] \gets A[i]\)
                \STATE \(i \gets i-1\)
            \ENDWHILE{}
            \STATE \(A[i + 1] \gets key\)
        \ENDFOR{}
        \RETURN \(A\)
    \end{algorithmic}
\end{algorithm}

الگوریتم~\ref{alg:inssort} نمونه‌ای از استفاده از یک الگوریتم را نشان می‌دهد.

\section{نتیجه‌گیری}
در این بخش نتیجه‌گیری فصل قرار می‌گیرد.